% Options for packages loaded elsewhere
% Options for packages loaded elsewhere
\PassOptionsToPackage{unicode}{hyperref}
\PassOptionsToPackage{hyphens}{url}
\PassOptionsToPackage{dvipsnames,svgnames,x11names}{xcolor}
%
\documentclass[
  letterpaper,
  DIV=11,
  numbers=noendperiod]{scrreprt}
\usepackage{xcolor}
\usepackage{amsmath,amssymb}
\setcounter{secnumdepth}{5}
\usepackage{iftex}
\ifPDFTeX
  \usepackage[T1]{fontenc}
  \usepackage[utf8]{inputenc}
  \usepackage{textcomp} % provide euro and other symbols
\else % if luatex or xetex
  \usepackage{unicode-math} % this also loads fontspec
  \defaultfontfeatures{Scale=MatchLowercase}
  \defaultfontfeatures[\rmfamily]{Ligatures=TeX,Scale=1}
\fi
\usepackage{lmodern}
\ifPDFTeX\else
  % xetex/luatex font selection
\fi
% Use upquote if available, for straight quotes in verbatim environments
\IfFileExists{upquote.sty}{\usepackage{upquote}}{}
\IfFileExists{microtype.sty}{% use microtype if available
  \usepackage[]{microtype}
  \UseMicrotypeSet[protrusion]{basicmath} % disable protrusion for tt fonts
}{}
\makeatletter
\@ifundefined{KOMAClassName}{% if non-KOMA class
  \IfFileExists{parskip.sty}{%
    \usepackage{parskip}
  }{% else
    \setlength{\parindent}{0pt}
    \setlength{\parskip}{6pt plus 2pt minus 1pt}}
}{% if KOMA class
  \KOMAoptions{parskip=half}}
\makeatother
% Make \paragraph and \subparagraph free-standing
\makeatletter
\ifx\paragraph\undefined\else
  \let\oldparagraph\paragraph
  \renewcommand{\paragraph}{
    \@ifstar
      \xxxParagraphStar
      \xxxParagraphNoStar
  }
  \newcommand{\xxxParagraphStar}[1]{\oldparagraph*{#1}\mbox{}}
  \newcommand{\xxxParagraphNoStar}[1]{\oldparagraph{#1}\mbox{}}
\fi
\ifx\subparagraph\undefined\else
  \let\oldsubparagraph\subparagraph
  \renewcommand{\subparagraph}{
    \@ifstar
      \xxxSubParagraphStar
      \xxxSubParagraphNoStar
  }
  \newcommand{\xxxSubParagraphStar}[1]{\oldsubparagraph*{#1}\mbox{}}
  \newcommand{\xxxSubParagraphNoStar}[1]{\oldsubparagraph{#1}\mbox{}}
\fi
\makeatother

\usepackage{color}
\usepackage{fancyvrb}
\newcommand{\VerbBar}{|}
\newcommand{\VERB}{\Verb[commandchars=\\\{\}]}
\DefineVerbatimEnvironment{Highlighting}{Verbatim}{commandchars=\\\{\}}
% Add ',fontsize=\small' for more characters per line
\usepackage{framed}
\definecolor{shadecolor}{RGB}{241,243,245}
\newenvironment{Shaded}{\begin{snugshade}}{\end{snugshade}}
\newcommand{\AlertTok}[1]{\textcolor[rgb]{0.68,0.00,0.00}{#1}}
\newcommand{\AnnotationTok}[1]{\textcolor[rgb]{0.37,0.37,0.37}{#1}}
\newcommand{\AttributeTok}[1]{\textcolor[rgb]{0.40,0.45,0.13}{#1}}
\newcommand{\BaseNTok}[1]{\textcolor[rgb]{0.68,0.00,0.00}{#1}}
\newcommand{\BuiltInTok}[1]{\textcolor[rgb]{0.00,0.23,0.31}{#1}}
\newcommand{\CharTok}[1]{\textcolor[rgb]{0.13,0.47,0.30}{#1}}
\newcommand{\CommentTok}[1]{\textcolor[rgb]{0.37,0.37,0.37}{#1}}
\newcommand{\CommentVarTok}[1]{\textcolor[rgb]{0.37,0.37,0.37}{\textit{#1}}}
\newcommand{\ConstantTok}[1]{\textcolor[rgb]{0.56,0.35,0.01}{#1}}
\newcommand{\ControlFlowTok}[1]{\textcolor[rgb]{0.00,0.23,0.31}{\textbf{#1}}}
\newcommand{\DataTypeTok}[1]{\textcolor[rgb]{0.68,0.00,0.00}{#1}}
\newcommand{\DecValTok}[1]{\textcolor[rgb]{0.68,0.00,0.00}{#1}}
\newcommand{\DocumentationTok}[1]{\textcolor[rgb]{0.37,0.37,0.37}{\textit{#1}}}
\newcommand{\ErrorTok}[1]{\textcolor[rgb]{0.68,0.00,0.00}{#1}}
\newcommand{\ExtensionTok}[1]{\textcolor[rgb]{0.00,0.23,0.31}{#1}}
\newcommand{\FloatTok}[1]{\textcolor[rgb]{0.68,0.00,0.00}{#1}}
\newcommand{\FunctionTok}[1]{\textcolor[rgb]{0.28,0.35,0.67}{#1}}
\newcommand{\ImportTok}[1]{\textcolor[rgb]{0.00,0.46,0.62}{#1}}
\newcommand{\InformationTok}[1]{\textcolor[rgb]{0.37,0.37,0.37}{#1}}
\newcommand{\KeywordTok}[1]{\textcolor[rgb]{0.00,0.23,0.31}{\textbf{#1}}}
\newcommand{\NormalTok}[1]{\textcolor[rgb]{0.00,0.23,0.31}{#1}}
\newcommand{\OperatorTok}[1]{\textcolor[rgb]{0.37,0.37,0.37}{#1}}
\newcommand{\OtherTok}[1]{\textcolor[rgb]{0.00,0.23,0.31}{#1}}
\newcommand{\PreprocessorTok}[1]{\textcolor[rgb]{0.68,0.00,0.00}{#1}}
\newcommand{\RegionMarkerTok}[1]{\textcolor[rgb]{0.00,0.23,0.31}{#1}}
\newcommand{\SpecialCharTok}[1]{\textcolor[rgb]{0.37,0.37,0.37}{#1}}
\newcommand{\SpecialStringTok}[1]{\textcolor[rgb]{0.13,0.47,0.30}{#1}}
\newcommand{\StringTok}[1]{\textcolor[rgb]{0.13,0.47,0.30}{#1}}
\newcommand{\VariableTok}[1]{\textcolor[rgb]{0.07,0.07,0.07}{#1}}
\newcommand{\VerbatimStringTok}[1]{\textcolor[rgb]{0.13,0.47,0.30}{#1}}
\newcommand{\WarningTok}[1]{\textcolor[rgb]{0.37,0.37,0.37}{\textit{#1}}}

\usepackage{longtable,booktabs,array}
\newcounter{none} % for unnumbered tables
\usepackage{calc} % for calculating minipage widths
% Correct order of tables after \paragraph or \subparagraph
\usepackage{etoolbox}
\makeatletter
\patchcmd\longtable{\par}{\if@noskipsec\mbox{}\fi\par}{}{}
\makeatother
% Allow footnotes in longtable head/foot
\IfFileExists{footnotehyper.sty}{\usepackage{footnotehyper}}{\usepackage{footnote}}
\makesavenoteenv{longtable}
\usepackage{graphicx}
\makeatletter
\newsavebox\pandoc@box
\newcommand*\pandocbounded[1]{% scales image to fit in text height/width
  \sbox\pandoc@box{#1}%
  \Gscale@div\@tempa{\textheight}{\dimexpr\ht\pandoc@box+\dp\pandoc@box\relax}%
  \Gscale@div\@tempb{\linewidth}{\wd\pandoc@box}%
  \ifdim\@tempb\p@<\@tempa\p@\let\@tempa\@tempb\fi% select the smaller of both
  \ifdim\@tempa\p@<\p@\scalebox{\@tempa}{\usebox\pandoc@box}%
  \else\usebox{\pandoc@box}%
  \fi%
}
% Set default figure placement to htbp
\def\fps@figure{htbp}
\makeatother





\setlength{\emergencystretch}{3em} % prevent overfull lines

\providecommand{\tightlist}{%
  \setlength{\itemsep}{0pt}\setlength{\parskip}{0pt}}



 


\KOMAoption{captions}{tableheading}
\makeatletter
\@ifpackageloaded{tcolorbox}{}{\usepackage[skins,breakable]{tcolorbox}}
\@ifpackageloaded{fontawesome5}{}{\usepackage{fontawesome5}}
\definecolor{quarto-callout-color}{HTML}{909090}
\definecolor{quarto-callout-note-color}{HTML}{0758E5}
\definecolor{quarto-callout-important-color}{HTML}{CC1914}
\definecolor{quarto-callout-warning-color}{HTML}{EB9113}
\definecolor{quarto-callout-tip-color}{HTML}{00A047}
\definecolor{quarto-callout-caution-color}{HTML}{FC5300}
\definecolor{quarto-callout-color-frame}{HTML}{acacac}
\definecolor{quarto-callout-note-color-frame}{HTML}{4582ec}
\definecolor{quarto-callout-important-color-frame}{HTML}{d9534f}
\definecolor{quarto-callout-warning-color-frame}{HTML}{f0ad4e}
\definecolor{quarto-callout-tip-color-frame}{HTML}{02b875}
\definecolor{quarto-callout-caution-color-frame}{HTML}{fd7e14}
\makeatother
\makeatletter
\@ifpackageloaded{bookmark}{}{\usepackage{bookmark}}
\makeatother
\makeatletter
\@ifpackageloaded{caption}{}{\usepackage{caption}}
\AtBeginDocument{%
\ifdefined\contentsname
  \renewcommand*\contentsname{Table of contents}
\else
  \newcommand\contentsname{Table of contents}
\fi
\ifdefined\listfigurename
  \renewcommand*\listfigurename{List of Figures}
\else
  \newcommand\listfigurename{List of Figures}
\fi
\ifdefined\listtablename
  \renewcommand*\listtablename{List of Tables}
\else
  \newcommand\listtablename{List of Tables}
\fi
\ifdefined\figurename
  \renewcommand*\figurename{Figure}
\else
  \newcommand\figurename{Figure}
\fi
\ifdefined\tablename
  \renewcommand*\tablename{Table}
\else
  \newcommand\tablename{Table}
\fi
}
\@ifpackageloaded{float}{}{\usepackage{float}}
\floatstyle{ruled}
\@ifundefined{c@chapter}{\newfloat{codelisting}{h}{lop}}{\newfloat{codelisting}{h}{lop}[chapter]}
\floatname{codelisting}{Listing}
\newcommand*\listoflistings{\listof{codelisting}{List of Listings}}
\makeatother
\makeatletter
\makeatother
\makeatletter
\@ifpackageloaded{caption}{}{\usepackage{caption}}
\@ifpackageloaded{subcaption}{}{\usepackage{subcaption}}
\makeatother
\usepackage{bookmark}
\IfFileExists{xurl.sty}{\usepackage{xurl}}{} % add URL line breaks if available
\urlstyle{same}
\hypersetup{
  pdftitle={Probabilidad I},
  pdfauthor={Matías Pozo},
  colorlinks=true,
  linkcolor={blue},
  filecolor={Maroon},
  citecolor={Blue},
  urlcolor={Blue},
  pdfcreator={LaTeX via pandoc}}


\title{Probabilidad I}
\author{Matías Pozo}
\date{2027-05-03}
\begin{document}
\maketitle

\renewcommand*\contentsname{Table of contents}
{
\setcounter{tocdepth}{2}
\tableofcontents
}

\bookmarksetup{startatroot}

\chapter*{Detalles del curso}\label{detalles-del-curso}
\addcontentsline{toc}{chapter}{Detalles del curso}

\markboth{Detalles del curso}{Detalles del curso}

Curso: Probabilidad I y II Siglas: EPG4101\\
Profesor: Johan Van Der Molen (jvandermolen@uc.cl)

\section*{Clases}\label{clases}
\addcontentsline{toc}{section}{Clases}

\markright{Clases}

\begin{itemize}
\tightlist
\item
  Clase 1 (I): Introducción (Lunes 16-03-2026)\\
\item
  Clase 2: Teoría de Conjuntos y Definición de Probabilidad (Lunes
  16-03-2026)
\item
  Clase 1 (II): Introducción (Lunes 25-05-2026)
\end{itemize}

\section*{Bibliografía}\label{bibliografuxeda}
\addcontentsline{toc}{section}{Bibliografía}

\markright{Bibliografía}

\begin{itemize}
\tightlist
\item
  Mínima:

  \begin{itemize}
  \tightlist
  \item
    Casella, G. y Berger, R.L. (2002). Statistical inference (2a.ed.).
    Cengage Learning.
    \href{https://drive.google.com/file/d/1ngxF_UTLxBLQcM0bQT6VJPwz60wVkMNV/view?usp=drive_link}{Libro}
  \end{itemize}
\item
  Complementaria:

  \begin{itemize}
  \tightlist
  \item
    Gut, A. (2009). An intermediate course in probability (2a.ed.)
    Springer.
    \href{https://drive.google.com/file/d/1QGeVrche29-WJN6JnBbaT295UWwxnvF-/view?usp=drive_link}{Libro}
  \item
    Ross, S.M. (2023). Simulation (6a.ed.). Academic Press.
    \href{https://drive.google.com/file/d/1ZG5jx-nEwqp8ssY_kc6HC-sW4qMfxbbj/view?usp=drive_link}{Libro}
  \end{itemize}
\end{itemize}

\bookmarksetup{startatroot}

\chapter{Introducción a la
probabilidad}\label{introducciuxf3n-a-la-probabilidad}

Existen distintas problemáticas o situaciones en el análisis de datos,
sin embargo, es posible generar preguntas generales para cualquier
problema que se presente:

\begin{itemize}
\tightlist
\item
  ¿Como abordar el problema?
\item
  ¿Existe una mejora real al aplicar un determinado factor en relación a
  uno existente o basal?
\item
  ¿Existe un cambio real del fenómeno actualmente observado, en relación
  a lo previamente conocido del fenómeno? (Considerando información
  previa y contrastandola con información actual que en un primer
  comienzo parece ser distinta a la habitual)
\item
  ¿Existe una anomalía?
\item
  ¿El tratamiento es realmente efectivo?
\end{itemize}

Todas estas preguntas conducen a si existe un real cambio, efecto,
mejora, disminución, anomalía\ldots{} Todas estas interpretaciones se
condensan en si existe o no una diferencia real.

En general, todos los problemas que enfrentaremos tienen algo en común,
y es que en todos vamos a necesitar observar datos y preguntarnos:

\begin{quote}
¿El resultado obtenido fue producto solo del azar o existe realmente un
fenómeno (natural o artificial) subyacente que explica los datos
observados?
\end{quote}

Para responder necesitamos lo siguiente:

\begin{itemize}
\tightlist
\item
  Definir un \textbf{modelo} para describir la variabilidad de los
  datos.
\item
  Una manera de \textbf{cuantificar} que los resultados observados son
  plausibles.
\item
  Una medida de \textbf{cuán sorprendente o raro} es lo observado.
\end{itemize}

\section{Teoría de la Probabilidad}\label{teoruxeda-de-la-probabilidad}

Para poder responder las preguntas que surjan de un problema o situación
en estudio, necesitamos la teoría de probabilidad, la cual nos entrega
las bases matemáticas para abordarlas.

La teoría de la probabilidad es el \textbf{lenguaje matemático de la
incertidumbre} y forman la base matemática de la inferencia estadística
y la ciencia de datos. Con la teoría de la probabilidad es posible:

\begin{itemize}
\tightlist
\item
  Modelar fenómenos aleatorios
\item
  Cuantificar incertidumbre. ¿Que tan seguros estamos de nuestra
  conclusión?
\item
  Entender la variabilidad en los datos
\item
  \textbf{Evaluar si un resultado observado es compatible con un modelo}
\end{itemize}

\begin{tcolorbox}[enhanced jigsaw, arc=.35mm, bottomrule=.15mm, bottomtitle=1mm, breakable, colback=white, colbacktitle=quarto-callout-note-color!10!white, colframe=quarto-callout-note-color-frame, coltitle=black, left=2mm, leftrule=.75mm, opacityback=0, opacitybacktitle=0.6, rightrule=.15mm, title=\textcolor{quarto-callout-note-color}{\faInfo}\hspace{0.5em}{Pensamiento importante}, titlerule=0mm, toprule=.15mm, toptitle=1mm]

Pude darme cuenta que es necesario preguntarse ¿Los grupos son
comparables? esto conlleva a pensar si la \textbf{comparación es justa}.
En este sentido, justa se refiere a que si es posible afirmar que lo
observado es \textbf{atribuible} solo al factor estudiado o no a otros
factores que podrián estar afectando a la variable de interes. Esto da
cuenta entonces de la necesidad de abordar un problema de manera en que
sea posible recabar información de variables importantes del problema y
que esas variables sean homogeneas en los grupos a comparar, esto
asegura que la comparación sea justa, que no se diferencien por otros
factores, y que entre ellos solo exista una diferencia en la variable o
variables que se busca estudiar y evaluar.

\end{tcolorbox}

Un aspecto muy importante es la elección de la distribución que se
utilizará para modelar el fenómeno aleatorio, ya que cada distribución
viene acompañada de supuestos que debemos considerar a la hora de
modelar el fenómeno.

\begin{tcolorbox}[enhanced jigsaw, arc=.35mm, bottomrule=.15mm, bottomtitle=1mm, breakable, colback=white, colbacktitle=quarto-callout-important-color!10!white, colframe=quarto-callout-important-color-frame, coltitle=black, left=2mm, leftrule=.75mm, opacityback=0, opacitybacktitle=0.6, rightrule=.15mm, title={Modelos, supuestos y simplificación}, titlerule=0mm, toprule=.15mm, toptitle=1mm]

Por lo tanto, como estadísticos proponemos un modelo matemático, el cual
intenta explicar un fenómeno. Dicho modelo siempre sera una
simplifiación del fenómeno, y cada supuesto expuesto suscribe
simplificación.

Por otro lado, es importante no caer en la sobresimplificación. Esto se
evita considerando los factores importantes del fenómeno que explican en
gran medida su variabilidad.

\end{tcolorbox}

\section{Historia de la probabilidad}\label{historia-de-la-probabilidad}

Los sumerios tenian juegos de \textbf{azar}. Ellos tenian un protodado,
el cual consistia en una especie de dado construido a partir de hueso de
animal marcado.\\
De hecho, la palabra azar proviene del árabe y significa flor, dado que
en esa época tallaban una flor en la cara que deseaban marcar.

Así, el origen de la teoría de la probabilidad está en los juegos de
azar.

En 1654, \textbf{Chevalier de Méré} comenzó a plantear interrogantes
acerca de cómo calcular probabilidades justas en los juegos de azar.

Luego, entre \textbf{Blaise Pascal} y \textbf{Pierre de Fermat}
comenzaron a analizar y desarrollar la teoría moderna de la
probabilidad.

El famoso \textbf{problema de los puntos} dio el punto de partida en
como abordar un problema, en el cual está presente el azar. Aquí Pascal
y Fermar propucieron distintas vías de resolución y ambas llegaron al
mismo resultado.

\begin{tcolorbox}[enhanced jigsaw, arc=.35mm, bottomrule=.15mm, bottomtitle=1mm, breakable, colback=white, colbacktitle=quarto-callout-tip-color!10!white, colframe=quarto-callout-tip-color-frame, coltitle=black, left=2mm, leftrule=.75mm, opacityback=0, opacitybacktitle=0.6, rightrule=.15mm, title=\textcolor{quarto-callout-tip-color}{\faLightbulb}\hspace{0.5em}{Problema de los puntos}, titlerule=0mm, toprule=.15mm, toptitle=1mm]

Dos jugadores apuestan dinero a un juego de azar (ej. lanzar una moneda)
con 5 rondas. El primero que gane 3 rondas se lleva todo el premio. Sin
embargo, el juego se interrumpe cuando el Jugador A ha ganado 2 rondas y
el Jugador B ha ganado 1. \textbf{¿como dividir la apuesta justamente si
el juego se detiene antes de finalizar?}

\end{tcolorbox}

En el año 1814, \textbf{Laplace} definió el concepto de probabilidad, la
cual corresponde a la razón entre \textbf{eventos favorable} y el total
de \textbf{eventos posibles}.

Laplace desarrolló:

\begin{itemize}
\tightlist
\item
  La Teoría de probabilidad inversas o inductivas (Teorema de Bayes)
\item
  Teorema Central del Límite
\item
  etc\ldots{}
\end{itemize}

\section{Interpretaciones de la probabilidad. Probabilidades objetivas y
subjetivas}\label{interpretaciones-de-la-probabilidad.-probabilidades-objetivas-y-subjetivas}

\begin{itemize}
\tightlist
\item
  Frecuencias relativas en infinitos experimentos
\item
  Grados de creencia
\end{itemize}

La base del pensamiento bayesiano es definir la probabilidad como un
\textbf{grado de creencia subjetivo} sobre un evento. Lo que busca es
actualizar estas estas creencias iniciales del fenómeno (\textbf{\emph{a
priori}}) mediante el uso de nueva evidencia (\textbf{\emph{datos}})
para obtener una creencia revisada (\textbf{\emph{a posteriori}}),
convirtiendo a la inferencia en un proceso dinámico y racional.

En términos un poco mas formales, se combina el conocimiento previo
(\(P(\Theta)\) o prior) con la verosimilitud de nuevos datos
(\(P(X \mid \Theta)\)) para generar la distribución posterior
(\(P(\Theta\mid X)\)).

En relación a cómo medir grados de creencia, Ramsey en el año 1926
definió probabilidad con un enfoque pragmático:

\begin{quote}
``Es la cantidad que se está dispuesto a pagar en una apuesta que paga
\$1''
\end{quote}

Además, demostro que para establecer coherencia, las probabilidades
asignadas por una misma persona deben cumplir con lo que conocemos hoy
en día como \textbf{Axiomas de probabilidad}.

Un aspecto interesante, en el caso del enfoque frecuentista, es que por
ejemplo, si consideramos una moneda con probabilidad 1/2, nunca
obtendremos una probabilidad 1/2 si tiramos la moneda una cantidad impar
de veces. Esto da cuenta de la naturaleza del fenómeno, dado que en la
realidad no puede darse. (Idea de experimento infinito con convergencia
de probabilidad de 1/2)

\section{Kolmogorov y los axiomas de
probabilidad}\label{kolmogorov-y-los-axiomas-de-probabilidad}

Hubo varios intentos para \textbf{axiomatizar la probabilidad}, hasta
que en 1933 \textbf{Kolmogorov publicó sus axiomas}, con una orientación
más abstracta, y se volvieron el nuevo estándar.

La axiomatización de la probabilidad permite conformar un cuerpo
matemático único, en donde toda interpretación de la probabilidad
(frecuentista, bayesiana, etc.) es irrelevantes, ya que los axiomas
brindan una consistencia para cualquiera de ellas.

Es decir, los axiomas de probabilidad de Kolmogorov establecen los
fundamentos matemáticos de la probabilidad. Los axiomas son 3:

\begin{tcolorbox}[enhanced jigsaw, arc=.35mm, bottomrule=.15mm, bottomtitle=1mm, breakable, colback=white, colbacktitle=quarto-callout-important-color!10!white, colframe=quarto-callout-important-color-frame, coltitle=black, left=2mm, leftrule=.75mm, opacityback=0, opacitybacktitle=0.6, rightrule=.15mm, title=\textcolor{quarto-callout-important-color}{\faExclamation}\hspace{0.5em}{Axiomas de Probabilidad de Kolmogorov Smirnov (1933)}, titlerule=0mm, toprule=.15mm, toptitle=1mm]

\begin{enumerate}
\def\labelenumi{\arabic{enumi}.}
\item
  \textbf{No negatividad:} La probabilidad de cualquier suceso \(A\) es
  no negativa: \(P(A) \geq 0\)
\item
  \textbf{Certeza:} La probabilidad del espacio muestral (suceso seguro)
  es 1: \(P(S) = 1\)
\item
  \textbf{Aditividad o \(\sigma-aditividad\):} Para sucesos mutuamente
  excluyentes, la probabilidad de su unión es la suma de sus
  probabilidades: \(P(A \cup B) = P(A) + P(B)\)\\
\end{enumerate}

\end{tcolorbox}

\bookmarksetup{startatroot}

\chapter{Teoría de Conjuntos y Definición de
Probabilidad}\label{teoruxeda-de-conjuntos-y-definiciuxf3n-de-probabilidad}

La teoría de probabilidad es la base sobre la cual está construida la
estadística. Esta nos brinda las herramientas para modelar una
población, un experimento o casi cualquier fenómeno aleatorio.

\begin{quote}
Teoría de probabilidad -\textgreater{} Base matemática para modelar un
fenómeno aleatorio.
\end{quote}

\begin{quote}
Mediante estos modelos, podemos ser capaces de extraer conclusiones
(inferencias) acerca de la población, inferencias basadas en la
observación de solo una parte del todo (una muestra de la población).
\end{quote}

Dado que la teoría de probabilidad está construida en base a la
\textbf{teoría de conjuntos}, sera necesario comenzar en ese punto, con
la finalidad de definir formalmente lo que llamamos \textbf{eventos} y
luego poder definir la \textbf{medida de probabilidad.}

\section{Teoría de Conjuntos}\label{teoruxeda-de-conjuntos}

\subsection{Conjunto}\label{conjunto}

\begin{tcolorbox}[enhanced jigsaw, arc=.35mm, bottomrule=.15mm, bottomtitle=1mm, breakable, colback=white, colbacktitle=quarto-callout-note-color!10!white, colframe=quarto-callout-note-color-frame, coltitle=black, left=2mm, leftrule=.75mm, opacityback=0, opacitybacktitle=0.6, rightrule=.15mm, title=\textcolor{quarto-callout-note-color}{\faInfo}\hspace{0.5em}{Definición de Conjunto}, titlerule=0mm, toprule=.15mm, toptitle=1mm]

Un conjunto es una colección definida de \textbf{objetos}. Los objetos
de un conjunto se denominan \textbf{elementos}.

\end{tcolorbox}

\begin{tcolorbox}[enhanced jigsaw, arc=.35mm, bottomrule=.15mm, bottomtitle=1mm, breakable, colback=white, colbacktitle=quarto-callout-warning-color!10!white, colframe=quarto-callout-warning-color-frame, coltitle=black, left=2mm, leftrule=.75mm, opacityback=0, opacitybacktitle=0.6, rightrule=.15mm, title={La probabilidad viene de los conjuntos}, titlerule=0mm, toprule=.15mm, toptitle=1mm]

De hecho, lo que hacemos es calcular un número a un conjunto, donde ese
número corresponde a la probabilidad asociada a ese conjunto.

\end{tcolorbox}

Si un objeto \(x\) pertenece al conjunto \(A\), se denota:

\[
x \in A
\] Y si no pertenece, se denota como:

\[
x \notin A
\] Un par de llaves \(\{ \}\), empleado junto con números, palabras o
símbolos, puede describir un conjunto. Por ejemplo:

\[
A = \{a_1, a_2, ... , a_n\}
\] También puede expresarse un conjunto \(A\) por medio de un símbolo
que denote una \textbf{variable}, en general la letra \(x\), la cual
representa a cualquier elemento de un conjunto y mediante la
\textbf{notación por construcción}, donde se emplea una barra vertical
\(|\) o \(:\), la cual siginifica \textbf{tales que:}.

Por ejemplo:

\[
S = \{x \ | \ x \ es \ un \ número \ natural \ menor \ que \ 6\}
\] Esto se lee así: \emph{El conjunto de todas las x tales que x es un
número natural menor que 6.}

Otro ejemplo:

\[
A = \{x \in \mathbb{N} : x \ es \ par \}
\]

\subsection{Subconjunto y conjunto de
partes}\label{subconjunto-y-conjunto-de-partes}

Sean \(A\) y \(B\) conjuntos, si todos los elementos de A están en B,
entonces se dice que \(A\) es \textbf{subconjunto} de \(B\), esto se
denota como:

\[
A \subseteq B \iff ( \forall x)(x \in A \Longrightarrow x \in B
)
\] Si \(A \subseteq B \ y \ A \neq  B\), entonces \(A\) se llama
\textbf{subconjunto propio} de B. También se denota \(A\subsetneq B\)

El \textbf{conjunto de partes} (o \textbf{conjunto potencia}) de \(A\),
denotado por \(\mathcal{P}(A)\), es el conjunto de \textbf{todos los
subconjuntos posibles} de \(A\), incluyendo el conjunto vacío
\((\emptyset)\) y el propio conjunto \(A\).

Si \(A\) tiene \(n\) elementos, entonces su conjunto de partes tiene
\(|\mathcal{P}(A)| = 2^n\) elementos. La cantidad de elementos del
conjunto de partes se conoce como \textbf{cardinalidad (\#)}.

Ejemplo: Si \(A = \{ 1, 2 \}\), entonces:\\
\[
\mathcal{P}(A) = \{ \emptyset, \{ 1\}, \{ 2\}, \{ 1,2\} \}
\] Tambien se lee a \(\mathcal{P}()\) como \textbf{las partes de}.

\section{Espacio Muestral y Evento}\label{espacio-muestral-y-evento}

Uno de los principales objetivos de un estadístico es poder extraer
conclusiones acerca de una población de \textbf{objetos} mediante la
realización de un experimento (obtención de una muestra de la
población).

El primer paso es poder identificar los \textbf{posibles resultados} de
un experimento, que en términos estadísticos se conoce como el
\textbf{espacio muestral}.

\begin{tcolorbox}[enhanced jigsaw, arc=.35mm, bottomrule=.15mm, bottomtitle=1mm, breakable, colback=white, colbacktitle=quarto-callout-note-color!10!white, colframe=quarto-callout-note-color-frame, coltitle=black, left=2mm, leftrule=.75mm, opacityback=0, opacitybacktitle=0.6, rightrule=.15mm, title={Espacio muestral}, titlerule=0mm, toprule=.15mm, toptitle=1mm]

El conjunto \(\Omega\) que contiene todos los resultados posibles de un
experimento particular se denomina \textbf{espacio muestral} de dicho
experimento.

\end{tcolorbox}

\begin{tcolorbox}[enhanced jigsaw, arc=.35mm, bottomrule=.15mm, bottomtitle=1mm, breakable, colback=white, colbacktitle=quarto-callout-note-color!10!white, colframe=quarto-callout-note-color-frame, coltitle=black, left=2mm, leftrule=.75mm, opacityback=0, opacitybacktitle=0.6, rightrule=.15mm, title={Evento}, titlerule=0mm, toprule=.15mm, toptitle=1mm]

Un evento es cualquier subconjunto \(A \subseteq \Omega\), es decir,
\(A \in \mathcal{P}(\Omega)\).

Podemos decir que el evento \(A\) ocurre si el resultado del experimento
o la muestra, está en el conjunto \(A\).

\end{tcolorbox}

Podemos clasificar espacios muestrales en función del número de
elementos que ellos contienen. Así, los espacios muestrales pueden ser
\textbf{contables (numerables)} o \textbf{no-contables (no-numerables)}.

\begin{tcolorbox}[enhanced jigsaw, arc=.35mm, bottomrule=.15mm, bottomtitle=1mm, breakable, colback=white, colbacktitle=quarto-callout-tip-color!10!white, colframe=quarto-callout-tip-color-frame, coltitle=black, left=2mm, leftrule=.75mm, opacityback=0, opacitybacktitle=0.6, rightrule=.15mm, title={Identificación de un conjunto contable}, titlerule=0mm, toprule=.15mm, toptitle=1mm]

Si los elementos de un espacio muestral pueden establecerse en una
correspondencia biunívoca con un subconjunto de los números naturales,
el espacio muestral es contable.

En otras palabras, si a cada elemento de \(\Omega\) es posible asignarle
un número natural como identificador, el conjunto es contable.

Obviamente, si un \(\Omega\) solo contiene un número finito de
elementos, el espacio muestral es contable.

\end{tcolorbox}

Esto implica que:\\
\$\$ \mathbb{N}, \mathbb{Z} ~y ~\mathbb{Q}, ~son ~~conjuntos ~contable
\textbackslash{} \mathbb{R} ~es ~No ~contable

\$\$

Saber si un espacio muestral es contable o no-contable nos servira para
entender como se dictan las maneras en que las probabilidades pueden ser
asignadas.

Hasta la actualidad, no es posible medir con precisión infinita una
cantidad física (Por ejemplo, el tiempo), por lo tanto, la mayor parte
del tiempo estamos trabajando con espacios muestrales contables.

Ocurre también que los cálculos para situaciones con mucha precisión, es
mejor considerarlo como un comportamiento continuo, mas que un conjunto
finito, dado que son muchos los valores que habria que considerar.

\section{Operaciones básicas con
eventos}\label{operaciones-buxe1sicas-con-eventos}

\begin{tcolorbox}[enhanced jigsaw, arc=.35mm, bottomrule=.15mm, bottomtitle=1mm, breakable, colback=white, colbacktitle=quarto-callout-note-color!10!white, colframe=quarto-callout-note-color-frame, coltitle=black, left=2mm, leftrule=.75mm, opacityback=0, opacitybacktitle=0.6, rightrule=.15mm, title={Igualdad}, titlerule=0mm, toprule=.15mm, toptitle=1mm]

Dados dos eventos \(A, B:\) \[
A = B \iff A \subseteq B \ \land B \subseteq A
\]

\end{tcolorbox}

\begin{tcolorbox}[enhanced jigsaw, arc=.35mm, bottomrule=.15mm, bottomtitle=1mm, breakable, colback=white, colbacktitle=quarto-callout-note-color!10!white, colframe=quarto-callout-note-color-frame, coltitle=black, left=2mm, leftrule=.75mm, opacityback=0, opacitybacktitle=0.6, rightrule=.15mm, title={Unión}, titlerule=0mm, toprule=.15mm, toptitle=1mm]

La unión de \(A\) y \(B\), que se escribe \(A \cup B\), es el conjunto
de elementos que pertenecen a \(A\), o a \(B\), o a ambos. \[
A \cup B = \{\omega \in \Omega : \omega \in A \ \lor \omega \in B \}
\]

\end{tcolorbox}

\begin{tcolorbox}[enhanced jigsaw, arc=.35mm, bottomrule=.15mm, bottomtitle=1mm, breakable, colback=white, colbacktitle=quarto-callout-note-color!10!white, colframe=quarto-callout-note-color-frame, coltitle=black, left=2mm, leftrule=.75mm, opacityback=0, opacitybacktitle=0.6, rightrule=.15mm, title={Intersección}, titlerule=0mm, toprule=.15mm, toptitle=1mm]

La intersección de \(A\) y \(B\), que se escribe \(A \cap B\), es el
conjunto de elementos que pertenecen tanto a \(A\) como a \(B\):

\[
A \cap B = \{\omega \in \Omega : \omega \in A \ \land \omega \in B \}
\]

\end{tcolorbox}

\begin{tcolorbox}[enhanced jigsaw, arc=.35mm, bottomrule=.15mm, bottomtitle=1mm, breakable, colback=white, colbacktitle=quarto-callout-note-color!10!white, colframe=quarto-callout-note-color-frame, coltitle=black, left=2mm, leftrule=.75mm, opacityback=0, opacitybacktitle=0.6, rightrule=.15mm, title={Complemento}, titlerule=0mm, toprule=.15mm, toptitle=1mm]

El complemento de \(A\), que se escribe \(A^c\), es el conjunto de todos
los elementos que no están en \(A\):

\[
A^c = \Omega \ \diagdown \ A = \{\omega \in \Omega : \omega \notin A\}
\]

\end{tcolorbox}

\begin{tcolorbox}[enhanced jigsaw, arc=.35mm, bottomrule=.15mm, bottomtitle=1mm, breakable, colback=white, colbacktitle=quarto-callout-note-color!10!white, colframe=quarto-callout-note-color-frame, coltitle=black, left=2mm, leftrule=.75mm, opacityback=0, opacitybacktitle=0.6, rightrule=.15mm, title={Diferencia}, titlerule=0mm, toprule=.15mm, toptitle=1mm]

La diferencia entre dos conjuntos \(A\) y \(B\) (denotada como
\(A \ \diagdown \ B\) es el nuevo conjunto formado por todos los
elementos que pertenecen al conjunto \(A\), pero \textbf{no} pertenecen
al conjunto \(B\).

\[
A \ \diagdown \ B = A \ \cap B^c
\]

\end{tcolorbox}

\section{Identidades y leyes
útiles}\label{identidades-y-leyes-uxfatiles}

\begin{tcolorbox}[enhanced jigsaw, arc=.35mm, bottomrule=.15mm, bottomtitle=1mm, breakable, colback=white, colbacktitle=quarto-callout-note-color!10!white, colframe=quarto-callout-note-color-frame, coltitle=black, left=2mm, leftrule=.75mm, opacityback=0, opacitybacktitle=0.6, rightrule=.15mm, title={Leyes de De Morgan}, titlerule=0mm, toprule=.15mm, toptitle=1mm]

\[(\cup_iA_i)^c = \cap_iA_i^c\] \[(\cap_iA_i)^c=\cup_iA_i^c\]

\end{tcolorbox}

\begin{tcolorbox}[enhanced jigsaw, arc=.35mm, bottomrule=.15mm, bottomtitle=1mm, breakable, colback=white, colbacktitle=quarto-callout-note-color!10!white, colframe=quarto-callout-note-color-frame, coltitle=black, left=2mm, leftrule=.75mm, opacityback=0, opacitybacktitle=0.6, rightrule=.15mm, title={Conmutatividad}, titlerule=0mm, toprule=.15mm, toptitle=1mm]

\[A \cup B = B \cup A\] \[A \cap B = B \cap A\]

\end{tcolorbox}

\begin{tcolorbox}[enhanced jigsaw, arc=.35mm, bottomrule=.15mm, bottomtitle=1mm, breakable, colback=white, colbacktitle=quarto-callout-note-color!10!white, colframe=quarto-callout-note-color-frame, coltitle=black, left=2mm, leftrule=.75mm, opacityback=0, opacitybacktitle=0.6, rightrule=.15mm, title={Asociatividad}, titlerule=0mm, toprule=.15mm, toptitle=1mm]

\[(A \cup B) \cup C = A \cup (B \cup C)\]
\[(A \cap B) \cap C = A \cap (B \cap C)\]

\end{tcolorbox}

\begin{tcolorbox}[enhanced jigsaw, arc=.35mm, bottomrule=.15mm, bottomtitle=1mm, breakable, colback=white, colbacktitle=quarto-callout-note-color!10!white, colframe=quarto-callout-note-color-frame, coltitle=black, left=2mm, leftrule=.75mm, opacityback=0, opacitybacktitle=0.6, rightrule=.15mm, title={Distributividad}, titlerule=0mm, toprule=.15mm, toptitle=1mm]

\[A \cup  (B \cap C) = (A \cup B) \cap (A \cup C)\]
\[A \cap  (B \cup C) = (A \cap B) \cup (A \cap C)\]

\end{tcolorbox}

\begin{tcolorbox}[enhanced jigsaw, arc=.35mm, bottomrule=.15mm, bottomtitle=1mm, breakable, colback=white, colbacktitle=quarto-callout-note-color!10!white, colframe=quarto-callout-note-color-frame, coltitle=black, left=2mm, leftrule=.75mm, opacityback=0, opacitybacktitle=0.6, rightrule=.15mm, title={Subconjunto de complementos}, titlerule=0mm, toprule=.15mm, toptitle=1mm]

Si \(A \subseteq B, entonces:\) \[
B^c \subseteq A^c
\]

\end{tcolorbox}

\section{Uniones e Intersecciones
infinitas}\label{uniones-e-intersecciones-infinitas}

Las operaciones de unión e intersección se pueden extender a una
colección infinita de conjuntos. Si \(A_1, A_2, A_3, \ldots\) es una
colección de conjuntos, todos definidas en un espacio muestral \(S\),
entonces existe una notación para una familia indexada
\(\{A_i\}_{i \in I} \ (índice \ I \ arbitrario):\)

\[
\bigcup_{i \in I}A_i = \{ x: \exists i \in I, x \in A_i \}
\] \[
\bigcap_{i \in I}A_i = \{ x: \forall  i \in I, x \in A_i \}
\] Para el caso contable (\textbf{es muy usual}): \[
\bigcup_{n=1}^{\infty} A_n \  \ y \ \ \bigcap_{n=1}^{\infty}A_n
\] También es posible aplicar la propiedad distributiva a una familia
infinita, la cual es válida para cualquier índice:\\
\[
Para \ todo \ conjunto \ B,
\] \[
B \cap(\bigcup_{i \in I} A_i) = \bigcup_{i \in I}(B \cap A_i) \ , \ \ \ \  B \cup(\bigcap_{i \in I} A_i) = \bigcap_{i \in I}(B \cup A_i)
\] \textbf{Las leyes de De Morgan también son válidas para familias
arbitrarias.}

Los ejemplos quedan pendientes, sin embargo para este tipo de
situaciones, la mayor parte del tiempo se emplean una familia de
intervalos como objetos matemáticos.

\section{Conjuntos disjuntos}\label{conjuntos-disjuntos}

\begin{tcolorbox}[enhanced jigsaw, arc=.35mm, bottomrule=.15mm, bottomtitle=1mm, breakable, colback=white, colbacktitle=quarto-callout-caution-color!10!white, colframe=quarto-callout-caution-color-frame, coltitle=black, left=2mm, leftrule=.75mm, opacityback=0, opacitybacktitle=0.6, rightrule=.15mm, title={Definición}, titlerule=0mm, toprule=.15mm, toptitle=1mm]

Dos conjuntos \(A\) y \(B\) son \textbf{mutuamente excluyentes} (o
\textbf{disjuntos}) si:\\
\[
A \cap B = \emptyset
\] Equivalente: No existe \(x\) tal que \(x \in A\) y \(x \in B\)
simultáneamente.

\end{tcolorbox}

Los eventos \(A_1, A_2, \ldots\) (\emph{que también se puede expresar
como una colección de conjuntos \(\{A_i\}_{i \in I}\)}) son
\textbf{disjuntos por pares} si cada par de posibles conjuntos distintos
tiene una intersección vacia, lo que significa que no comparten ningún
elemento en común.

En notación matemática, los eventos disjuntos por pares se representan
como: \[
A_i \cap A_j = \emptyset, \ \forall i \neq j
\]

\section{\texorpdfstring{Partición de un conjunto (Partición de
\(\Omega\) en el caso de considerar al espacio
muestral)}{Partición de un conjunto (Partición de \textbackslash Omega en el caso de considerar al espacio muestral)}}\label{particiuxf3n-de-un-conjunto-particiuxf3n-de-omega-en-el-caso-de-considerar-al-espacio-muestral}

A partir del concepto de eventos disjuntos es posible construir el
concepto de partición de un conjunto.

\begin{tcolorbox}[enhanced jigsaw, arc=.35mm, bottomrule=.15mm, bottomtitle=1mm, breakable, colback=white, colbacktitle=quarto-callout-warning-color!10!white, colframe=quarto-callout-warning-color-frame, coltitle=black, left=2mm, leftrule=.75mm, opacityback=0, opacitybacktitle=0.6, rightrule=.15mm, title={Definición de Partición de un conjunto}, titlerule=0mm, toprule=.15mm, toptitle=1mm]

Una colección \(\{A_i\}_{i \in I}\) de subconjuntos de \(\omega\) es una
\textbf{partición} de \(\omega\) si se cumplen \textbf{tres
condiciones:}

\begin{enumerate}
\def\labelenumi{\arabic{enumi}.}
\tightlist
\item
  \(A_i = \emptyset, \ \forall i\) (Ningún subconjunto es vacio).
\item
  \(A_i \cap A_j = \emptyset, \ \forall i \neq j\) (Son subconjuntos
  disjuntos por pares, mutuamente excluyentes).
\item
  \(\bigcup_{i \in I}A_i = \Omega\) (La unión de los subconjuntos forma
  el espacio muestral).
\end{enumerate}

\end{tcolorbox}

Siempre es posible hacer una partición, cuando no existe, por medio de
partir en secciones más finas. La noción intuitiva seria pensar en un
rompecabezas.

\begin{tcolorbox}[enhanced jigsaw, arc=.35mm, bottomrule=.15mm, bottomtitle=1mm, breakable, colback=white, colbacktitle=quarto-callout-tip-color!10!white, colframe=quarto-callout-tip-color-frame, coltitle=black, left=2mm, leftrule=.75mm, opacityback=0, opacitybacktitle=0.6, rightrule=.15mm, title=\textcolor{quarto-callout-tip-color}{\faLightbulb}\hspace{0.5em}{Ejemplos de particiones en \(\mathbb{R}\)}, titlerule=0mm, toprule=.15mm, toptitle=1mm]

Las siguientes familias o colecciones de intervalos son particiones de
\(\mathbb{R}\): \[
\{(-\infty,0), \{0\}, (0, +\infty)\}
\] \[
A_i = [i,\  i+1], \ \ i \in \mathbb{Z}
\]

\end{tcolorbox}

\section{Sigma-álgebra}\label{sigma-uxe1lgebra}

El objeto matemático \textbf{sigma-álgebra es primordial en la
definición de probabilidad}.

\begin{tcolorbox}[enhanced jigsaw, arc=.35mm, bottomrule=.15mm, bottomtitle=1mm, breakable, colback=white, colbacktitle=quarto-callout-warning-color!10!white, colframe=quarto-callout-warning-color-frame, coltitle=black, left=2mm, leftrule=.75mm, opacityback=0, opacitybacktitle=0.6, rightrule=.15mm, title={Definición de Sigma-álgebra}, titlerule=0mm, toprule=.15mm, toptitle=1mm]

Sea \(\Omega \neq \emptyset\) un espacio muestral. Decimos que
\(\mathcal{F}\) es una \textbf{Sigma-álgebra de partes} de \(\Omega\)
(es decir, \(\mathcal{F}\) es una colección de partes de \(\Omega\),
\textbf{\(\mathcal{F}\) es un conjunto de subconjunto de las partes de
\(\Omega\)}, contiene subconjuntos de \(\Omega\)) si \textbf{cumple tres
condiciones:}

\begin{enumerate}
\def\labelenumi{\arabic{enumi}.}
\tightlist
\item
  \(\Omega \in \mathcal{F}\) (Contiene al espacio muestral).
\item
  \(A \in \mathcal{F} \implies A^c \in \mathcal{F}\) (Cerrada para
  complementos).
\item
  \(\{A_i\}_{i \in I} \in \mathcal{F} \implies \bigcup_{i = 1}^{\infty}A_i \in \mathcal{F}\)
  (Cerrada para uniones infinitas numerables/contables).
\end{enumerate}

\textbf{Observaciones:}

\begin{itemize}
\tightlist
\item
  \(\emptyset\) pertenece a cualquier Sigma-álgebra, dado que el punto 2
  indica que el complemento de cualquier conjunto debe estar, en este
  caso el complemento de \(\Omega\) es \(\emptyset\)
\item
  Por las leyes de De Morgan, las Sigma-álgebras \textbf{también son
  cerradas para intersecciones numerables}. Esto es por que como las
  uniones tienen que estar, los complementos tienen que estar, entonces,
  los complementos de las uniones tambien tienen que estar, es decir,
  las intersecciones.
\item
  Para un mismo \(\Omega\) pueden existir muchas Sigma-álgebra.
\item
  \(\mathcal{P}(\Omega)\) es siempre una Sigma-álgebra.
\end{itemize}

\end{tcolorbox}

\begin{tcolorbox}[enhanced jigsaw, arc=.35mm, bottomrule=.15mm, bottomtitle=1mm, breakable, colback=white, colbacktitle=quarto-callout-tip-color!10!white, colframe=quarto-callout-tip-color-frame, coltitle=black, left=2mm, leftrule=.75mm, opacityback=0, opacitybacktitle=0.6, rightrule=.15mm, title={Ejemplos de elección de Sigma-álgebra}, titlerule=0mm, toprule=.15mm, toptitle=1mm]

\begin{itemize}
\tightlist
\item
  Si \(\Omega\) es finito o contable, entonces típicamente escogemos la
  sigma-álgebra de las partes.
\item
  Si \(\Omega = \mathbb{R}\), entonces escogemos todas las uniones e
  intersecciones de conjuntos de la forma intervalar:
  \([a, b], \ (a, b], \ [a, b), \ (a, b])\)
\end{itemize}

\end{tcolorbox}

\begin{tcolorbox}[enhanced jigsaw, arc=.35mm, bottomrule=.15mm, bottomtitle=1mm, breakable, colback=white, colbacktitle=quarto-callout-important-color!10!white, colframe=quarto-callout-important-color-frame, coltitle=black, left=2mm, leftrule=.75mm, opacityback=0, opacitybacktitle=0.6, rightrule=.15mm, title=\textcolor{quarto-callout-important-color}{\faExclamation}\hspace{0.5em}{Sigma-álgebra como base para las probabilidades}, titlerule=0mm, toprule=.15mm, toptitle=1mm]

\begin{quote}
El objetivo de definir una Sigma-álgebra \(\mathcal{F}\) de las partes
de \(\Omega\) es para poder calcular probabilidades.
\end{quote}

\begin{itemize}
\tightlist
\item
  \textbf{La Sigma-álgebra va a ser mi conjunto de todos los conjuntos a
  los que yo les puedo calcular una probabilidad.}
\item
  \textbf{La Sigma-álgebra contiene a todos los conjuntos a los que yo
  le voy a calcular la probabilidad.}
\item
  \textbf{Es el conjunto de todos los elementos que me interesa calcular
  una probabilidad.}
\item
  \textbf{La Sigma-álgebra es el conjunto de todos los eventos a los
  cuales quiero calcular la probabilidad.}
\end{itemize}

\end{tcolorbox}

\section{Axiomas de Probabilidad}\label{axiomas-de-probabilidad}

\begin{tcolorbox}[enhanced jigsaw, arc=.35mm, bottomrule=.15mm, bottomtitle=1mm, breakable, colback=white, colbacktitle=quarto-callout-warning-color!10!white, colframe=quarto-callout-warning-color-frame, coltitle=black, left=2mm, leftrule=.75mm, opacityback=0, opacitybacktitle=0.6, rightrule=.15mm, title={Definición}, titlerule=0mm, toprule=.15mm, toptitle=1mm]

Una medida de probabilidad \(P\) en un espacio muestral \(\Omega\) es
una función definida en una Sigma-álgebra
\(\mathcal{F} \subseteq \mathcal{P(\Omega)}\) tal que:

\begin{enumerate}
\def\labelenumi{\arabic{enumi}.}
\tightlist
\item
  \(P(A) \ge 0, \ \forall A \in \mathcal{F}\)
\item
  \(P(\Omega) = 1\)
\item
  Si \(A_1, A_2,\ldots\) es una secuencia de eventos mutuamente
  excluyentes en \(\mathcal{F}\), entonces,
  \(P(\bigcup_{i = 1}^{\infty}A_i)=\sum_{i=1}^{\infty}P(A_i)\)
\end{enumerate}

El trío \(( \Omega, \mathcal{F} \ y \ P)\) de objetos matemáticos se
denominan \textbf{espacio de probabilidad.}

\end{tcolorbox}

\section{Método general para definir
probabilidades}\label{muxe9todo-general-para-definir-probabilidades}

Para poder definir una medida de probabilidad a un espacio muestral
basta con asignar un valor a cada elemento del espacio muestral, y hay
que regirse por los axiomas de probabilidad, es decir, todos deben ser
positivos y la suma de todos debe ser 1.

\begin{tcolorbox}[enhanced jigsaw, arc=.35mm, bottomrule=.15mm, bottomtitle=1mm, breakable, colback=white, colbacktitle=quarto-callout-warning-color!10!white, colframe=quarto-callout-warning-color-frame, coltitle=black, left=2mm, leftrule=.75mm, opacityback=0, opacitybacktitle=0.6, rightrule=.15mm, title={Teorema}, titlerule=0mm, toprule=.15mm, toptitle=1mm]

Sea \(\Omega = \{s_1, \ldots,s_n\}\) un conjunto finito. Sea
\(\mathcal{F}\) cualquier Sigma-álgebra de partes de \(\Omega\). Sea
además \(p_1, p_2, \ldots, p_n\) números positivos, tal que
\(\sum_{i=1}^np_i=1\). Para cada \(A \in \mathcal{F}\) definimos:

\[
P(A) = \sum_{i:s_i \in A }p_i
\] Entonces \(P\) es una medida de probabilidad. Además, el resultado
también es válido para espacios muestrales contables, ya que contable no
implica finito.

\end{tcolorbox}

\bookmarksetup{startatroot}

\chapter{Cálculo de Probabilidades}\label{cuxe1lculo-de-probabilidades}

\section{Cálculo de probabilidades en el enfoque
axiomático}\label{cuxe1lculo-de-probabilidades-en-el-enfoque-axiomuxe1tico}

\begin{itemize}
\tightlist
\item
  La probabilidad axiomática no solo define una función \(P\), también
  nos \textbf{brinda un conjunto de reglas y propiedades} para calcular
  probabilidades complejas.
\item
  El conjunto de reglas se organizan en base a \textbf{identidades
  elementales sobre conjuntos} y de \textbf{consecuencias directas de
  los axiomas.}
\item
  El objetivo es poder \textbf{transformar un sistema de axiomas en una
  herramienta de cálculo y aproximación.}
\end{itemize}

\section{Notación y marco general}\label{notaciuxf3n-y-marco-general}

\bookmarksetup{startatroot}

\chapter{clase\_1\_II}\label{clase_1_ii}

\bookmarksetup{startatroot}

\chapter{Introducción}\label{introducciuxf3n}

Considerando la infinidad de problemas que se pueden abordar en el
contexto del análisis de datos, lo importante a tener siempre en
consideración son los siguientes aspectos:

\begin{enumerate}
\def\labelenumi{\arabic{enumi}.}
\tightlist
\item
  Que variables vas a observar
\item
  Como vas a relacionar estas variables
\end{enumerate}

\bookmarksetup{startatroot}

\chapter{Summary}\label{summary}

In summary, this book has no content whatsoever.

\begin{Shaded}
\begin{Highlighting}[]
\DecValTok{1} \SpecialCharTok{+} \DecValTok{1}
\end{Highlighting}
\end{Shaded}

\begin{verbatim}
[1] 2
\end{verbatim}

\bookmarksetup{startatroot}

\chapter*{References}\label{references}
\addcontentsline{toc}{chapter}{References}

\markboth{References}{References}

\protect\phantomsection\label{refs}




\end{document}
